\documentclass[tikz,border=0pt]{standalone}
\usepackage[T1]{fontenc}
\usepackage{lmodern}
% mqttsystem-figure-system.tex
% Canonical TikZ visual system for MQTTSuite technical figures.
%
% Usage:
%   \documentclass[tikz,border=10pt]{standalone}
%   % mqttsystem-figure-system.tex
% Canonical TikZ visual system for MQTTSuite technical figures.
%
% Usage:
%   \documentclass[tikz,border=10pt]{standalone}
%   % mqttsystem-figure-system.tex
% Canonical TikZ visual system for MQTTSuite technical figures.
%
% Usage:
%   \documentclass[tikz,border=10pt]{standalone}
%   \input{mqttsystem-figure-system.tex}
%   \begin{document}
%   \begin{tikzpicture}[mqtt desktop]
%     ...
%   \end{tikzpicture}
%   \end{document}
%
% Canonical source rule:
% - figure .tex sources + this file are authoritative;
% - generated SVGs are publication output and are never hand-edited;
% - real runtime/UI evidence remains real raster capture.

\usetikzlibrary{
  arrows.meta,
  backgrounds,
  calc,
  fit,
  matrix,
  positioning,
  shapes.geometric,
  shapes.misc,
  shadows.blur
}

% ---------------------------------------------------------------------------
% Canonical palette
% ---------------------------------------------------------------------------

\definecolor{MQText}{HTML}{0F172A}
\definecolor{MQTextMuted}{HTML}{475569}
\definecolor{MQLine}{HTML}{475569}
\definecolor{MQBorder}{HTML}{CBD5E1}

\definecolor{MQSurface}{HTML}{F8FAFC}
\definecolor{MQSurfaceAlt}{HTML}{F1F5F9}
\definecolor{MQWhite}{HTML}{FFFFFF}

\definecolor{MQDecisionFill}{HTML}{FFF7ED}
\definecolor{MQDecisionLine}{HTML}{D97706}

\definecolor{MQSuccessFill}{HTML}{ECFDF5}
\definecolor{MQSuccessLine}{HTML}{10B981}

\definecolor{MQErrorFill}{HTML}{FEF2F2}
\definecolor{MQErrorLine}{HTML}{EF4444}

\definecolor{MQControlFill}{HTML}{EFF6FF}
\definecolor{MQControlLine}{HTML}{3B82F6}

\definecolor{MQWarningFill}{HTML}{FFF7ED}
\definecolor{MQWarningLine}{HTML}{F59E0B}

\definecolor{MQCrossLane}{HTML}{2563EB}

\definecolor{MQDataPlaneFill}{HTML}{EFF6FF}
\definecolor{MQDataPlaneLine}{HTML}{2563EB}

\definecolor{MQAdminPlaneFill}{HTML}{F8FAFC}
\definecolor{MQAdminPlaneLine}{HTML}{64748B}

\definecolor{MQTrustFill}{HTML}{FFF7ED}
\definecolor{MQTrustLine}{HTML}{D97706}

\definecolor{MQEvidenceRuntimeFill}{HTML}{ECFDF5}
\definecolor{MQEvidenceRuntimeLine}{HTML}{10B981}

\definecolor{MQEvidenceSourceFill}{HTML}{EFF6FF}
\definecolor{MQEvidenceSourceLine}{HTML}{3B82F6}

% ---------------------------------------------------------------------------
% Canonical radius system
% ---------------------------------------------------------------------------
%
% Small informational elements stay subtle; ordinary nodes are moderately
% rounded; structural containers are more rounded. Decision diamonds and
% database cylinders do not use these radii.

\newcommand{\MQRadiusXS}{1.2pt}
\newcommand{\MQRadiusS}{1.8pt}
\newcommand{\MQRadiusM}{2.8pt}
\newcommand{\MQRadiusL}{5.0pt}
\newcommand{\MQRadiusXL}{7.0pt}

% ---------------------------------------------------------------------------
% Canonical shading / elevation system
% ---------------------------------------------------------------------------
%
% Flat rendering is the default. Shadows are opt-in and reserved for semantic
% emphasis, never general decoration.

\newcommand{\MQShadowXSubtle}{0.7pt}
\newcommand{\MQShadowYSubtle}{-0.7pt}
\newcommand{\MQShadowBlurSubtle}{1.6pt}
\newcommand{\MQShadowOpacitySubtle}{0.12}

\newcommand{\MQShadowXRaised}{1.0pt}
\newcommand{\MQShadowYRaised}{-1.0pt}
\newcommand{\MQShadowBlurRaised}{2.4pt}
\newcommand{\MQShadowOpacityRaised}{0.16}

\tikzset{
  mqtt elevated none/.style={},
  mqtt elevated subtle/.style={
    blur shadow={
      shadow xshift=\MQShadowXSubtle,
      shadow yshift=\MQShadowYSubtle,
      shadow blur radius=\MQShadowBlurSubtle,
      opacity=\MQShadowOpacitySubtle
    }
  },
  mqtt elevated raised/.style={
    blur shadow={
      shadow xshift=\MQShadowXRaised,
      shadow yshift=\MQShadowYRaised,
      shadow blur radius=\MQShadowBlurRaised,
      opacity=\MQShadowOpacityRaised
    }
  }
}

% ---------------------------------------------------------------------------
% Canonical spacing system
% ---------------------------------------------------------------------------
%
% Base grid: 1 spacing unit = 2 mm. Ordinary figure spacing should use these
% semantic constants rather than ad-hoc gaps wherever practical. Exact geometry
% remains valid when a figure needs deterministic alignment, but should still
% follow this grid.

\newcommand{\MQU}{2mm}

% Generic scale
\newcommand{\MQGapXS}{2\MQU}      % 4 mm
\newcommand{\MQGapS}{3\MQU}       % 6 mm
\newcommand{\MQGapM}{5\MQU}       % 10 mm
\newcommand{\MQGapL}{7\MQU}       % 14 mm
\newcommand{\MQGapXL}{10\MQU}     % 20 mm

% Node interior
\newcommand{\MQNodePadX}{2.5\MQU} % 5 mm
\newcommand{\MQNodePadY}{2\MQU}   % 4 mm

% Sequential node spacing
\newcommand{\MQNodeGapTight}{3\MQU} % 6 mm
\newcommand{\MQNodeGap}{5\MQU}      % 10 mm
\newcommand{\MQNodeGapLoose}{7\MQU} % 14 mm

% Decision / branch geometry
\newcommand{\MQBranchGap}{7\MQU}    % 14 mm
\newcommand{\MQBranchRun}{5\MQU}    % 10 mm orthogonal run before turn

% Lanes and semantic containers
\newcommand{\MQLaneGap}{12\MQU}     % 24 mm between major lanes
\newcommand{\MQLanePadX}{4\MQU}     % 8 mm
\newcommand{\MQLanePadY}{4\MQU}     % 8 mm
\newcommand{\MQGroupPad}{3\MQU}     % 6 mm
\newcommand{\MQTrustPad}{4\MQU}     % 8 mm

% Titles and labels
\newcommand{\MQTitleGap}{2\MQU}     % 4 mm
\newcommand{\MQLabelGap}{1\MQU}     % 2 mm
\newcommand{\MQCrossLaneGap}{4\MQU} % 8 mm reserved clearance

% Figure safe area
\newcommand{\MQFigureSafeX}{5\MQU}  % 10 mm
\newcommand{\MQFigureSafeY}{5\MQU}  % 10 mm

% Compatibility aliases
\newcommand{\MQDesktopSafe}{\MQFigureSafeX}
\newcommand{\MQMobileSafe}{4\MQU}   % 8 mm

% Mobile art direction: mobile is not a scaled desktop composition.
\newcommand{\MQMobileNodeGap}{6\MQU}   % 12 mm
\newcommand{\MQMobileBranchGap}{8\MQU} % 16 mm
\newcommand{\MQMobileLaneGap}{10\MQU}  % 20 mm
\newcommand{\MQMobileGroupPad}{4\MQU}  % 8 mm

% ---------------------------------------------------------------------------
% Typography
% ---------------------------------------------------------------------------

\tikzset{
  mqtt text/.style={
    font=\sffamily\small,
    text=MQText,
    align=center
  },
  mqtt text left/.style={
    mqtt text,
    align=left
  },
  mqtt title/.style={
    font=\sffamily\bfseries\large,
    text=MQText,
    align=left
  },
  mqtt section title/.style={
    font=\sffamily\bfseries\small,
    text=MQText,
    align=left,
    anchor=west
  },
  mqtt annotation/.style={
    font=\sffamily\scriptsize,
    text=MQTextMuted,
    align=center
  },
  mqtt code/.style={
    font=\ttfamily\small,
    text=MQText
  },
  mqtt code small/.style={
    font=\ttfamily\scriptsize,
    text=MQText
  }
}

% ---------------------------------------------------------------------------
% Figure presets
% ---------------------------------------------------------------------------
%
% These establish common defaults only. Desktop/mobile variants may and often
% should use different geometry.

\tikzset{
  mqtt desktop/.style={
    x=1cm,
    y=1cm,
    every node/.append style={transform shape=false}
  },
  mqtt mobile/.style={
    x=1cm,
    y=1cm,
    every node/.append style={
      transform shape=false,
      font=\sffamily\normalsize
    }
  }
}

% ---------------------------------------------------------------------------
% Core node families
% ---------------------------------------------------------------------------

\tikzset{
  mqtt process/.style={
    mqtt text,
    draw=MQLine,
    fill=MQWhite,
    rounded corners=\MQRadiusM,
    line width=.75pt,
    minimum height=10mm,
    minimum width=25mm,
    inner xsep=\MQNodePadX,
    inner ysep=\MQNodePadY
  },
  mqtt application/.style={
    mqtt process,
    minimum width=31mm,
    minimum height=13mm,
    line width=.9pt
  },
  mqtt external/.style={
    mqtt process,
    fill=MQSurfaceAlt,
    draw=MQBorder
  },
  mqtt state/.style={
    mqtt process,
    fill=MQSurface
  },
  mqtt decision/.style={
    mqtt text,
    diamond,
    aspect=1.75,
    draw=MQDecisionLine,
    fill=MQDecisionFill,
    line width=.8pt,
    inner sep=1.2pt,
    minimum width=27mm,
    minimum height=14mm
  },
  mqtt success/.style={
    mqtt process,
    fill=MQSuccessFill,
    draw=MQSuccessLine
  },
  mqtt error/.style={
    mqtt process,
    fill=MQErrorFill,
    draw=MQErrorLine
  },
  mqtt control/.style={
    mqtt process,
    fill=MQControlFill,
    draw=MQControlLine
  },
  mqtt warning/.style={
    mqtt process,
    fill=MQWarningFill,
    draw=MQWarningLine
  },
  mqtt note/.style={
    mqtt text,
    draw=MQBorder,
    fill=MQSurfaceAlt,
    rounded corners=\MQRadiusS,
    line width=.65pt,
    font=\sffamily\scriptsize,
    minimum height=8mm,
    inner xsep=\MQNodePadX,
    inner ysep=\MQNodePadY
  },
  mqtt config/.style={
    mqtt process,
    fill=MQSurfaceAlt,
    draw=MQBorder,
    rounded corners=\MQRadiusS
  },
  mqtt data/.style={
    mqtt process,
    fill=MQWhite,
    draw=MQBorder
  },
  mqtt database/.style={
    mqtt text,
    cylinder,
    shape border rotate=90,
    aspect=.25,
    draw=MQLine,
    fill=MQWhite,
    line width=.75pt,
    minimum width=24mm,
    minimum height=13mm,
    inner xsep=\MQNodePadX,
    inner ysep=\MQNodePadY
  }
}

% ---------------------------------------------------------------------------
% Semantic containers
% ---------------------------------------------------------------------------

\tikzset{
  mqtt lane/.style={
    draw=MQBorder,
    fill=MQSurface,
    rounded corners=\MQRadiusL,
    line width=.8pt,
    inner xsep=\MQLanePadX,
    inner ysep=\MQLanePadY
  },
  mqtt group/.style={
    draw=MQBorder,
    fill=MQWhite,
    rounded corners=\MQRadiusL,
    line width=.75pt,
    inner sep=\MQGroupPad
  },
  mqtt data plane/.style={
    draw=MQDataPlaneLine,
    fill=MQDataPlaneFill,
    rounded corners=\MQRadiusL,
    line width=.9pt,
    inner sep=\MQGroupPad
  },
  mqtt admin plane/.style={
    draw=MQAdminPlaneLine,
    fill=MQAdminPlaneFill,
    rounded corners=\MQRadiusL,
    line width=.9pt,
    inner sep=\MQGroupPad
  },
  mqtt trust boundary/.style={
    draw=MQTrustLine,
    fill=none,
    rounded corners=\MQRadiusXL,
    line width=1.0pt,
    dashed,
    inner sep=\MQTrustPad
  }
}

% ---------------------------------------------------------------------------
% Arrow semantics
% ---------------------------------------------------------------------------

\tikzset{
  mqtt flow/.style={
    -{Stealth[length=2.2mm,width=1.5mm]},
    draw=MQLine,
    line width=.8pt
  },
  mqtt data flow/.style={
    -{Stealth[length=2.2mm,width=1.5mm]},
    draw=MQDataPlaneLine,
    line width=.9pt
  },
  mqtt control flow/.style={
    -{Stealth[length=2.2mm,width=1.5mm]},
    draw=MQControlLine,
    line width=.9pt
  },
  mqtt observation/.style={
    -{Stealth[length=2.0mm,width=1.4mm]},
    draw=MQLine,
    line width=.75pt,
    dashed
  },
  mqtt handoff/.style={
    -{Stealth[length=2.3mm,width=1.6mm]},
    draw=MQCrossLane,
    line width=1.0pt
  },
  mqtt association/.style={
    draw=MQBorder,
    line width=.7pt
  },
  mqtt dependency/.style={
    -{Stealth[length=2.0mm,width=1.4mm]},
    draw=MQLine,
    line width=.7pt,
    dashed
  }
}

% ---------------------------------------------------------------------------
% Labels and evidence
% ---------------------------------------------------------------------------

\tikzset{
  mqtt branch label/.style={
    font=\sffamily\scriptsize\bfseries,
    text=MQTextMuted,
    fill=MQWhite,
    inner sep=1.4pt
  },
  mqtt edge label/.style={
    font=\sffamily\scriptsize,
    text=MQTextMuted,
    fill=MQWhite,
    inner sep=1.4pt
  },
  mqtt topic label/.style={
    font=\ttfamily\scriptsize,
    text=MQText,
    fill=MQWhite,
    inner sep=1.4pt
  },
  mqtt route label/.style={
    font=\ttfamily\scriptsize,
    text=MQText,
    fill=MQWhite,
    inner sep=1.4pt
  },
  mqtt evidence runtime/.style={
    font=\sffamily\scriptsize\bfseries,
    text=MQSuccessLine,
    draw=MQEvidenceRuntimeLine,
    fill=MQEvidenceRuntimeFill,
    rounded corners=\MQRadiusXS,
    line width=.6pt,
    inner xsep=4pt,
    inner ysep=2pt
  },
  mqtt evidence source/.style={
    font=\sffamily\scriptsize\bfseries,
    text=MQControlLine,
    draw=MQEvidenceSourceLine,
    fill=MQEvidenceSourceFill,
    rounded corners=\MQRadiusXS,
    line width=.6pt,
    inner xsep=4pt,
    inner ysep=2pt
  }
}

% ---------------------------------------------------------------------------
% Canonical publication rules
% ---------------------------------------------------------------------------
%
% 1. Figure .tex + this style file are canonical technical sources.
% 2. Generated SVGs are never hand-edited.
% 3. SVG is the publication format for technical diagrams.
% 4. PNG is reserved for genuine runtime/UI evidence or raster-only input.
% 5. Desktop/mobile may share semantics but use independently art-directed geometry.
% 6. No arrow exists without real directional semantics.
% 7. Dashed arrows mean observation/secondary relation, not primary flow.
% 8. Containers denote real process/runtime/trust/configuration/lifecycle boundaries.
% 9. Color never carries the only semantic distinction.
% 10. Data plane, admin/control plane, trust boundaries and evidence classes use
%     the canonical semantic styles above.
% 11. Prefer explicit anchors/ports and orthogonal routing where clarity improves.
% 12. Preserve readable text size at actual GitHub desktop/mobile rendering widths.
% 13. Ordinary spacing follows the semantic spacing system above; avoid ad-hoc gaps.
% 14. Rounded corners follow the semantic radius scale above; avoid one-off radii.
% 15. Flat rendering is the default; elevation/shadows are explicit and semantically
%     justified.

%   \begin{document}
%   \begin{tikzpicture}[mqtt desktop]
%     ...
%   \end{tikzpicture}
%   \end{document}
%
% Canonical source rule:
% - figure .tex sources + this file are authoritative;
% - generated SVGs are publication output and are never hand-edited;
% - real runtime/UI evidence remains real raster capture.

\usetikzlibrary{
  arrows.meta,
  backgrounds,
  calc,
  fit,
  matrix,
  positioning,
  shapes.geometric,
  shapes.misc,
  shadows.blur
}

% ---------------------------------------------------------------------------
% Canonical palette
% ---------------------------------------------------------------------------

\definecolor{MQText}{HTML}{0F172A}
\definecolor{MQTextMuted}{HTML}{475569}
\definecolor{MQLine}{HTML}{475569}
\definecolor{MQBorder}{HTML}{CBD5E1}

\definecolor{MQSurface}{HTML}{F8FAFC}
\definecolor{MQSurfaceAlt}{HTML}{F1F5F9}
\definecolor{MQWhite}{HTML}{FFFFFF}

\definecolor{MQDecisionFill}{HTML}{FFF7ED}
\definecolor{MQDecisionLine}{HTML}{D97706}

\definecolor{MQSuccessFill}{HTML}{ECFDF5}
\definecolor{MQSuccessLine}{HTML}{10B981}

\definecolor{MQErrorFill}{HTML}{FEF2F2}
\definecolor{MQErrorLine}{HTML}{EF4444}

\definecolor{MQControlFill}{HTML}{EFF6FF}
\definecolor{MQControlLine}{HTML}{3B82F6}

\definecolor{MQWarningFill}{HTML}{FFF7ED}
\definecolor{MQWarningLine}{HTML}{F59E0B}

\definecolor{MQCrossLane}{HTML}{2563EB}

\definecolor{MQDataPlaneFill}{HTML}{EFF6FF}
\definecolor{MQDataPlaneLine}{HTML}{2563EB}

\definecolor{MQAdminPlaneFill}{HTML}{F8FAFC}
\definecolor{MQAdminPlaneLine}{HTML}{64748B}

\definecolor{MQTrustFill}{HTML}{FFF7ED}
\definecolor{MQTrustLine}{HTML}{D97706}

\definecolor{MQEvidenceRuntimeFill}{HTML}{ECFDF5}
\definecolor{MQEvidenceRuntimeLine}{HTML}{10B981}

\definecolor{MQEvidenceSourceFill}{HTML}{EFF6FF}
\definecolor{MQEvidenceSourceLine}{HTML}{3B82F6}

% ---------------------------------------------------------------------------
% Canonical radius system
% ---------------------------------------------------------------------------
%
% Small informational elements stay subtle; ordinary nodes are moderately
% rounded; structural containers are more rounded. Decision diamonds and
% database cylinders do not use these radii.

\newcommand{\MQRadiusXS}{1.2pt}
\newcommand{\MQRadiusS}{1.8pt}
\newcommand{\MQRadiusM}{2.8pt}
\newcommand{\MQRadiusL}{5.0pt}
\newcommand{\MQRadiusXL}{7.0pt}

% ---------------------------------------------------------------------------
% Canonical shading / elevation system
% ---------------------------------------------------------------------------
%
% Flat rendering is the default. Shadows are opt-in and reserved for semantic
% emphasis, never general decoration.

\newcommand{\MQShadowXSubtle}{0.7pt}
\newcommand{\MQShadowYSubtle}{-0.7pt}
\newcommand{\MQShadowBlurSubtle}{1.6pt}
\newcommand{\MQShadowOpacitySubtle}{0.12}

\newcommand{\MQShadowXRaised}{1.0pt}
\newcommand{\MQShadowYRaised}{-1.0pt}
\newcommand{\MQShadowBlurRaised}{2.4pt}
\newcommand{\MQShadowOpacityRaised}{0.16}

\tikzset{
  mqtt elevated none/.style={},
  mqtt elevated subtle/.style={
    blur shadow={
      shadow xshift=\MQShadowXSubtle,
      shadow yshift=\MQShadowYSubtle,
      shadow blur radius=\MQShadowBlurSubtle,
      opacity=\MQShadowOpacitySubtle
    }
  },
  mqtt elevated raised/.style={
    blur shadow={
      shadow xshift=\MQShadowXRaised,
      shadow yshift=\MQShadowYRaised,
      shadow blur radius=\MQShadowBlurRaised,
      opacity=\MQShadowOpacityRaised
    }
  }
}

% ---------------------------------------------------------------------------
% Canonical spacing system
% ---------------------------------------------------------------------------
%
% Base grid: 1 spacing unit = 2 mm. Ordinary figure spacing should use these
% semantic constants rather than ad-hoc gaps wherever practical. Exact geometry
% remains valid when a figure needs deterministic alignment, but should still
% follow this grid.

\newcommand{\MQU}{2mm}

% Generic scale
\newcommand{\MQGapXS}{2\MQU}      % 4 mm
\newcommand{\MQGapS}{3\MQU}       % 6 mm
\newcommand{\MQGapM}{5\MQU}       % 10 mm
\newcommand{\MQGapL}{7\MQU}       % 14 mm
\newcommand{\MQGapXL}{10\MQU}     % 20 mm

% Node interior
\newcommand{\MQNodePadX}{2.5\MQU} % 5 mm
\newcommand{\MQNodePadY}{2\MQU}   % 4 mm

% Sequential node spacing
\newcommand{\MQNodeGapTight}{3\MQU} % 6 mm
\newcommand{\MQNodeGap}{5\MQU}      % 10 mm
\newcommand{\MQNodeGapLoose}{7\MQU} % 14 mm

% Decision / branch geometry
\newcommand{\MQBranchGap}{7\MQU}    % 14 mm
\newcommand{\MQBranchRun}{5\MQU}    % 10 mm orthogonal run before turn

% Lanes and semantic containers
\newcommand{\MQLaneGap}{12\MQU}     % 24 mm between major lanes
\newcommand{\MQLanePadX}{4\MQU}     % 8 mm
\newcommand{\MQLanePadY}{4\MQU}     % 8 mm
\newcommand{\MQGroupPad}{3\MQU}     % 6 mm
\newcommand{\MQTrustPad}{4\MQU}     % 8 mm

% Titles and labels
\newcommand{\MQTitleGap}{2\MQU}     % 4 mm
\newcommand{\MQLabelGap}{1\MQU}     % 2 mm
\newcommand{\MQCrossLaneGap}{4\MQU} % 8 mm reserved clearance

% Figure safe area
\newcommand{\MQFigureSafeX}{5\MQU}  % 10 mm
\newcommand{\MQFigureSafeY}{5\MQU}  % 10 mm

% Compatibility aliases
\newcommand{\MQDesktopSafe}{\MQFigureSafeX}
\newcommand{\MQMobileSafe}{4\MQU}   % 8 mm

% Mobile art direction: mobile is not a scaled desktop composition.
\newcommand{\MQMobileNodeGap}{6\MQU}   % 12 mm
\newcommand{\MQMobileBranchGap}{8\MQU} % 16 mm
\newcommand{\MQMobileLaneGap}{10\MQU}  % 20 mm
\newcommand{\MQMobileGroupPad}{4\MQU}  % 8 mm

% ---------------------------------------------------------------------------
% Typography
% ---------------------------------------------------------------------------

\tikzset{
  mqtt text/.style={
    font=\sffamily\small,
    text=MQText,
    align=center
  },
  mqtt text left/.style={
    mqtt text,
    align=left
  },
  mqtt title/.style={
    font=\sffamily\bfseries\large,
    text=MQText,
    align=left
  },
  mqtt section title/.style={
    font=\sffamily\bfseries\small,
    text=MQText,
    align=left,
    anchor=west
  },
  mqtt annotation/.style={
    font=\sffamily\scriptsize,
    text=MQTextMuted,
    align=center
  },
  mqtt code/.style={
    font=\ttfamily\small,
    text=MQText
  },
  mqtt code small/.style={
    font=\ttfamily\scriptsize,
    text=MQText
  }
}

% ---------------------------------------------------------------------------
% Figure presets
% ---------------------------------------------------------------------------
%
% These establish common defaults only. Desktop/mobile variants may and often
% should use different geometry.

\tikzset{
  mqtt desktop/.style={
    x=1cm,
    y=1cm,
    every node/.append style={transform shape=false}
  },
  mqtt mobile/.style={
    x=1cm,
    y=1cm,
    every node/.append style={
      transform shape=false,
      font=\sffamily\normalsize
    }
  }
}

% ---------------------------------------------------------------------------
% Core node families
% ---------------------------------------------------------------------------

\tikzset{
  mqtt process/.style={
    mqtt text,
    draw=MQLine,
    fill=MQWhite,
    rounded corners=\MQRadiusM,
    line width=.75pt,
    minimum height=10mm,
    minimum width=25mm,
    inner xsep=\MQNodePadX,
    inner ysep=\MQNodePadY
  },
  mqtt application/.style={
    mqtt process,
    minimum width=31mm,
    minimum height=13mm,
    line width=.9pt
  },
  mqtt external/.style={
    mqtt process,
    fill=MQSurfaceAlt,
    draw=MQBorder
  },
  mqtt state/.style={
    mqtt process,
    fill=MQSurface
  },
  mqtt decision/.style={
    mqtt text,
    diamond,
    aspect=1.75,
    draw=MQDecisionLine,
    fill=MQDecisionFill,
    line width=.8pt,
    inner sep=1.2pt,
    minimum width=27mm,
    minimum height=14mm
  },
  mqtt success/.style={
    mqtt process,
    fill=MQSuccessFill,
    draw=MQSuccessLine
  },
  mqtt error/.style={
    mqtt process,
    fill=MQErrorFill,
    draw=MQErrorLine
  },
  mqtt control/.style={
    mqtt process,
    fill=MQControlFill,
    draw=MQControlLine
  },
  mqtt warning/.style={
    mqtt process,
    fill=MQWarningFill,
    draw=MQWarningLine
  },
  mqtt note/.style={
    mqtt text,
    draw=MQBorder,
    fill=MQSurfaceAlt,
    rounded corners=\MQRadiusS,
    line width=.65pt,
    font=\sffamily\scriptsize,
    minimum height=8mm,
    inner xsep=\MQNodePadX,
    inner ysep=\MQNodePadY
  },
  mqtt config/.style={
    mqtt process,
    fill=MQSurfaceAlt,
    draw=MQBorder,
    rounded corners=\MQRadiusS
  },
  mqtt data/.style={
    mqtt process,
    fill=MQWhite,
    draw=MQBorder
  },
  mqtt database/.style={
    mqtt text,
    cylinder,
    shape border rotate=90,
    aspect=.25,
    draw=MQLine,
    fill=MQWhite,
    line width=.75pt,
    minimum width=24mm,
    minimum height=13mm,
    inner xsep=\MQNodePadX,
    inner ysep=\MQNodePadY
  }
}

% ---------------------------------------------------------------------------
% Semantic containers
% ---------------------------------------------------------------------------

\tikzset{
  mqtt lane/.style={
    draw=MQBorder,
    fill=MQSurface,
    rounded corners=\MQRadiusL,
    line width=.8pt,
    inner xsep=\MQLanePadX,
    inner ysep=\MQLanePadY
  },
  mqtt group/.style={
    draw=MQBorder,
    fill=MQWhite,
    rounded corners=\MQRadiusL,
    line width=.75pt,
    inner sep=\MQGroupPad
  },
  mqtt data plane/.style={
    draw=MQDataPlaneLine,
    fill=MQDataPlaneFill,
    rounded corners=\MQRadiusL,
    line width=.9pt,
    inner sep=\MQGroupPad
  },
  mqtt admin plane/.style={
    draw=MQAdminPlaneLine,
    fill=MQAdminPlaneFill,
    rounded corners=\MQRadiusL,
    line width=.9pt,
    inner sep=\MQGroupPad
  },
  mqtt trust boundary/.style={
    draw=MQTrustLine,
    fill=none,
    rounded corners=\MQRadiusXL,
    line width=1.0pt,
    dashed,
    inner sep=\MQTrustPad
  }
}

% ---------------------------------------------------------------------------
% Arrow semantics
% ---------------------------------------------------------------------------

\tikzset{
  mqtt flow/.style={
    -{Stealth[length=2.2mm,width=1.5mm]},
    draw=MQLine,
    line width=.8pt
  },
  mqtt data flow/.style={
    -{Stealth[length=2.2mm,width=1.5mm]},
    draw=MQDataPlaneLine,
    line width=.9pt
  },
  mqtt control flow/.style={
    -{Stealth[length=2.2mm,width=1.5mm]},
    draw=MQControlLine,
    line width=.9pt
  },
  mqtt observation/.style={
    -{Stealth[length=2.0mm,width=1.4mm]},
    draw=MQLine,
    line width=.75pt,
    dashed
  },
  mqtt handoff/.style={
    -{Stealth[length=2.3mm,width=1.6mm]},
    draw=MQCrossLane,
    line width=1.0pt
  },
  mqtt association/.style={
    draw=MQBorder,
    line width=.7pt
  },
  mqtt dependency/.style={
    -{Stealth[length=2.0mm,width=1.4mm]},
    draw=MQLine,
    line width=.7pt,
    dashed
  }
}

% ---------------------------------------------------------------------------
% Labels and evidence
% ---------------------------------------------------------------------------

\tikzset{
  mqtt branch label/.style={
    font=\sffamily\scriptsize\bfseries,
    text=MQTextMuted,
    fill=MQWhite,
    inner sep=1.4pt
  },
  mqtt edge label/.style={
    font=\sffamily\scriptsize,
    text=MQTextMuted,
    fill=MQWhite,
    inner sep=1.4pt
  },
  mqtt topic label/.style={
    font=\ttfamily\scriptsize,
    text=MQText,
    fill=MQWhite,
    inner sep=1.4pt
  },
  mqtt route label/.style={
    font=\ttfamily\scriptsize,
    text=MQText,
    fill=MQWhite,
    inner sep=1.4pt
  },
  mqtt evidence runtime/.style={
    font=\sffamily\scriptsize\bfseries,
    text=MQSuccessLine,
    draw=MQEvidenceRuntimeLine,
    fill=MQEvidenceRuntimeFill,
    rounded corners=\MQRadiusXS,
    line width=.6pt,
    inner xsep=4pt,
    inner ysep=2pt
  },
  mqtt evidence source/.style={
    font=\sffamily\scriptsize\bfseries,
    text=MQControlLine,
    draw=MQEvidenceSourceLine,
    fill=MQEvidenceSourceFill,
    rounded corners=\MQRadiusXS,
    line width=.6pt,
    inner xsep=4pt,
    inner ysep=2pt
  }
}

% ---------------------------------------------------------------------------
% Canonical publication rules
% ---------------------------------------------------------------------------
%
% 1. Figure .tex + this style file are canonical technical sources.
% 2. Generated SVGs are never hand-edited.
% 3. SVG is the publication format for technical diagrams.
% 4. PNG is reserved for genuine runtime/UI evidence or raster-only input.
% 5. Desktop/mobile may share semantics but use independently art-directed geometry.
% 6. No arrow exists without real directional semantics.
% 7. Dashed arrows mean observation/secondary relation, not primary flow.
% 8. Containers denote real process/runtime/trust/configuration/lifecycle boundaries.
% 9. Color never carries the only semantic distinction.
% 10. Data plane, admin/control plane, trust boundaries and evidence classes use
%     the canonical semantic styles above.
% 11. Prefer explicit anchors/ports and orthogonal routing where clarity improves.
% 12. Preserve readable text size at actual GitHub desktop/mobile rendering widths.
% 13. Ordinary spacing follows the semantic spacing system above; avoid ad-hoc gaps.
% 14. Rounded corners follow the semantic radius scale above; avoid one-off radii.
% 15. Flat rendering is the default; elevation/shadows are explicit and semantically
%     justified.

%   \begin{document}
%   \begin{tikzpicture}[mqtt desktop]
%     ...
%   \end{tikzpicture}
%   \end{document}
%
% Canonical source rule:
% - figure .tex sources + this file are authoritative;
% - generated SVGs are publication output and are never hand-edited;
% - real runtime/UI evidence remains real raster capture.

\usetikzlibrary{
  arrows.meta,
  backgrounds,
  calc,
  fit,
  matrix,
  positioning,
  shapes.geometric,
  shapes.misc,
  shadows.blur
}

% ---------------------------------------------------------------------------
% Canonical palette
% ---------------------------------------------------------------------------

\definecolor{MQText}{HTML}{0F172A}
\definecolor{MQTextMuted}{HTML}{475569}
\definecolor{MQLine}{HTML}{475569}
\definecolor{MQBorder}{HTML}{CBD5E1}

\definecolor{MQSurface}{HTML}{F8FAFC}
\definecolor{MQSurfaceAlt}{HTML}{F1F5F9}
\definecolor{MQWhite}{HTML}{FFFFFF}

\definecolor{MQDecisionFill}{HTML}{FFF7ED}
\definecolor{MQDecisionLine}{HTML}{D97706}

\definecolor{MQSuccessFill}{HTML}{ECFDF5}
\definecolor{MQSuccessLine}{HTML}{10B981}

\definecolor{MQErrorFill}{HTML}{FEF2F2}
\definecolor{MQErrorLine}{HTML}{EF4444}

\definecolor{MQControlFill}{HTML}{EFF6FF}
\definecolor{MQControlLine}{HTML}{3B82F6}

\definecolor{MQWarningFill}{HTML}{FFF7ED}
\definecolor{MQWarningLine}{HTML}{F59E0B}

\definecolor{MQCrossLane}{HTML}{2563EB}

\definecolor{MQDataPlaneFill}{HTML}{EFF6FF}
\definecolor{MQDataPlaneLine}{HTML}{2563EB}

\definecolor{MQAdminPlaneFill}{HTML}{F8FAFC}
\definecolor{MQAdminPlaneLine}{HTML}{64748B}

\definecolor{MQTrustFill}{HTML}{FFF7ED}
\definecolor{MQTrustLine}{HTML}{D97706}

\definecolor{MQEvidenceRuntimeFill}{HTML}{ECFDF5}
\definecolor{MQEvidenceRuntimeLine}{HTML}{10B981}

\definecolor{MQEvidenceSourceFill}{HTML}{EFF6FF}
\definecolor{MQEvidenceSourceLine}{HTML}{3B82F6}

% ---------------------------------------------------------------------------
% Canonical radius system
% ---------------------------------------------------------------------------
%
% Small informational elements stay subtle; ordinary nodes are moderately
% rounded; structural containers are more rounded. Decision diamonds and
% database cylinders do not use these radii.

\newcommand{\MQRadiusXS}{1.2pt}
\newcommand{\MQRadiusS}{1.8pt}
\newcommand{\MQRadiusM}{2.8pt}
\newcommand{\MQRadiusL}{5.0pt}
\newcommand{\MQRadiusXL}{7.0pt}

% ---------------------------------------------------------------------------
% Canonical shading / elevation system
% ---------------------------------------------------------------------------
%
% Flat rendering is the default. Shadows are opt-in and reserved for semantic
% emphasis, never general decoration.

\newcommand{\MQShadowXSubtle}{0.7pt}
\newcommand{\MQShadowYSubtle}{-0.7pt}
\newcommand{\MQShadowBlurSubtle}{1.6pt}
\newcommand{\MQShadowOpacitySubtle}{0.12}

\newcommand{\MQShadowXRaised}{1.0pt}
\newcommand{\MQShadowYRaised}{-1.0pt}
\newcommand{\MQShadowBlurRaised}{2.4pt}
\newcommand{\MQShadowOpacityRaised}{0.16}

\tikzset{
  mqtt elevated none/.style={},
  mqtt elevated subtle/.style={
    blur shadow={
      shadow xshift=\MQShadowXSubtle,
      shadow yshift=\MQShadowYSubtle,
      shadow blur radius=\MQShadowBlurSubtle,
      opacity=\MQShadowOpacitySubtle
    }
  },
  mqtt elevated raised/.style={
    blur shadow={
      shadow xshift=\MQShadowXRaised,
      shadow yshift=\MQShadowYRaised,
      shadow blur radius=\MQShadowBlurRaised,
      opacity=\MQShadowOpacityRaised
    }
  }
}

% ---------------------------------------------------------------------------
% Canonical spacing system
% ---------------------------------------------------------------------------
%
% Base grid: 1 spacing unit = 2 mm. Ordinary figure spacing should use these
% semantic constants rather than ad-hoc gaps wherever practical. Exact geometry
% remains valid when a figure needs deterministic alignment, but should still
% follow this grid.

\newcommand{\MQU}{2mm}

% Generic scale
\newcommand{\MQGapXS}{2\MQU}      % 4 mm
\newcommand{\MQGapS}{3\MQU}       % 6 mm
\newcommand{\MQGapM}{5\MQU}       % 10 mm
\newcommand{\MQGapL}{7\MQU}       % 14 mm
\newcommand{\MQGapXL}{10\MQU}     % 20 mm

% Node interior
\newcommand{\MQNodePadX}{2.5\MQU} % 5 mm
\newcommand{\MQNodePadY}{2\MQU}   % 4 mm

% Sequential node spacing
\newcommand{\MQNodeGapTight}{3\MQU} % 6 mm
\newcommand{\MQNodeGap}{5\MQU}      % 10 mm
\newcommand{\MQNodeGapLoose}{7\MQU} % 14 mm

% Decision / branch geometry
\newcommand{\MQBranchGap}{7\MQU}    % 14 mm
\newcommand{\MQBranchRun}{5\MQU}    % 10 mm orthogonal run before turn

% Lanes and semantic containers
\newcommand{\MQLaneGap}{12\MQU}     % 24 mm between major lanes
\newcommand{\MQLanePadX}{4\MQU}     % 8 mm
\newcommand{\MQLanePadY}{4\MQU}     % 8 mm
\newcommand{\MQGroupPad}{3\MQU}     % 6 mm
\newcommand{\MQTrustPad}{4\MQU}     % 8 mm

% Titles and labels
\newcommand{\MQTitleGap}{2\MQU}     % 4 mm
\newcommand{\MQLabelGap}{1\MQU}     % 2 mm
\newcommand{\MQCrossLaneGap}{4\MQU} % 8 mm reserved clearance

% Figure safe area
\newcommand{\MQFigureSafeX}{5\MQU}  % 10 mm
\newcommand{\MQFigureSafeY}{5\MQU}  % 10 mm

% Compatibility aliases
\newcommand{\MQDesktopSafe}{\MQFigureSafeX}
\newcommand{\MQMobileSafe}{4\MQU}   % 8 mm

% Mobile art direction: mobile is not a scaled desktop composition.
\newcommand{\MQMobileNodeGap}{6\MQU}   % 12 mm
\newcommand{\MQMobileBranchGap}{8\MQU} % 16 mm
\newcommand{\MQMobileLaneGap}{10\MQU}  % 20 mm
\newcommand{\MQMobileGroupPad}{4\MQU}  % 8 mm

% ---------------------------------------------------------------------------
% Typography
% ---------------------------------------------------------------------------

\tikzset{
  mqtt text/.style={
    font=\sffamily\small,
    text=MQText,
    align=center
  },
  mqtt text left/.style={
    mqtt text,
    align=left
  },
  mqtt title/.style={
    font=\sffamily\bfseries\large,
    text=MQText,
    align=left
  },
  mqtt section title/.style={
    font=\sffamily\bfseries\small,
    text=MQText,
    align=left,
    anchor=west
  },
  mqtt annotation/.style={
    font=\sffamily\scriptsize,
    text=MQTextMuted,
    align=center
  },
  mqtt code/.style={
    font=\ttfamily\small,
    text=MQText
  },
  mqtt code small/.style={
    font=\ttfamily\scriptsize,
    text=MQText
  }
}

% ---------------------------------------------------------------------------
% Figure presets
% ---------------------------------------------------------------------------
%
% These establish common defaults only. Desktop/mobile variants may and often
% should use different geometry.

\tikzset{
  mqtt desktop/.style={
    x=1cm,
    y=1cm,
    every node/.append style={transform shape=false}
  },
  mqtt mobile/.style={
    x=1cm,
    y=1cm,
    every node/.append style={
      transform shape=false,
      font=\sffamily\normalsize
    }
  }
}

% ---------------------------------------------------------------------------
% Core node families
% ---------------------------------------------------------------------------

\tikzset{
  mqtt process/.style={
    mqtt text,
    draw=MQLine,
    fill=MQWhite,
    rounded corners=\MQRadiusM,
    line width=.75pt,
    minimum height=10mm,
    minimum width=25mm,
    inner xsep=\MQNodePadX,
    inner ysep=\MQNodePadY
  },
  mqtt application/.style={
    mqtt process,
    minimum width=31mm,
    minimum height=13mm,
    line width=.9pt
  },
  mqtt external/.style={
    mqtt process,
    fill=MQSurfaceAlt,
    draw=MQBorder
  },
  mqtt state/.style={
    mqtt process,
    fill=MQSurface
  },
  mqtt decision/.style={
    mqtt text,
    diamond,
    aspect=1.75,
    draw=MQDecisionLine,
    fill=MQDecisionFill,
    line width=.8pt,
    inner sep=1.2pt,
    minimum width=27mm,
    minimum height=14mm
  },
  mqtt success/.style={
    mqtt process,
    fill=MQSuccessFill,
    draw=MQSuccessLine
  },
  mqtt error/.style={
    mqtt process,
    fill=MQErrorFill,
    draw=MQErrorLine
  },
  mqtt control/.style={
    mqtt process,
    fill=MQControlFill,
    draw=MQControlLine
  },
  mqtt warning/.style={
    mqtt process,
    fill=MQWarningFill,
    draw=MQWarningLine
  },
  mqtt note/.style={
    mqtt text,
    draw=MQBorder,
    fill=MQSurfaceAlt,
    rounded corners=\MQRadiusS,
    line width=.65pt,
    font=\sffamily\scriptsize,
    minimum height=8mm,
    inner xsep=\MQNodePadX,
    inner ysep=\MQNodePadY
  },
  mqtt config/.style={
    mqtt process,
    fill=MQSurfaceAlt,
    draw=MQBorder,
    rounded corners=\MQRadiusS
  },
  mqtt data/.style={
    mqtt process,
    fill=MQWhite,
    draw=MQBorder
  },
  mqtt database/.style={
    mqtt text,
    cylinder,
    shape border rotate=90,
    aspect=.25,
    draw=MQLine,
    fill=MQWhite,
    line width=.75pt,
    minimum width=24mm,
    minimum height=13mm,
    inner xsep=\MQNodePadX,
    inner ysep=\MQNodePadY
  }
}

% ---------------------------------------------------------------------------
% Semantic containers
% ---------------------------------------------------------------------------

\tikzset{
  mqtt lane/.style={
    draw=MQBorder,
    fill=MQSurface,
    rounded corners=\MQRadiusL,
    line width=.8pt,
    inner xsep=\MQLanePadX,
    inner ysep=\MQLanePadY
  },
  mqtt group/.style={
    draw=MQBorder,
    fill=MQWhite,
    rounded corners=\MQRadiusL,
    line width=.75pt,
    inner sep=\MQGroupPad
  },
  mqtt data plane/.style={
    draw=MQDataPlaneLine,
    fill=MQDataPlaneFill,
    rounded corners=\MQRadiusL,
    line width=.9pt,
    inner sep=\MQGroupPad
  },
  mqtt admin plane/.style={
    draw=MQAdminPlaneLine,
    fill=MQAdminPlaneFill,
    rounded corners=\MQRadiusL,
    line width=.9pt,
    inner sep=\MQGroupPad
  },
  mqtt trust boundary/.style={
    draw=MQTrustLine,
    fill=none,
    rounded corners=\MQRadiusXL,
    line width=1.0pt,
    dashed,
    inner sep=\MQTrustPad
  }
}

% ---------------------------------------------------------------------------
% Arrow semantics
% ---------------------------------------------------------------------------

\tikzset{
  mqtt flow/.style={
    -{Stealth[length=2.2mm,width=1.5mm]},
    draw=MQLine,
    line width=.8pt
  },
  mqtt data flow/.style={
    -{Stealth[length=2.2mm,width=1.5mm]},
    draw=MQDataPlaneLine,
    line width=.9pt
  },
  mqtt control flow/.style={
    -{Stealth[length=2.2mm,width=1.5mm]},
    draw=MQControlLine,
    line width=.9pt
  },
  mqtt observation/.style={
    -{Stealth[length=2.0mm,width=1.4mm]},
    draw=MQLine,
    line width=.75pt,
    dashed
  },
  mqtt handoff/.style={
    -{Stealth[length=2.3mm,width=1.6mm]},
    draw=MQCrossLane,
    line width=1.0pt
  },
  mqtt association/.style={
    draw=MQBorder,
    line width=.7pt
  },
  mqtt dependency/.style={
    -{Stealth[length=2.0mm,width=1.4mm]},
    draw=MQLine,
    line width=.7pt,
    dashed
  }
}

% ---------------------------------------------------------------------------
% Labels and evidence
% ---------------------------------------------------------------------------

\tikzset{
  mqtt branch label/.style={
    font=\sffamily\scriptsize\bfseries,
    text=MQTextMuted,
    fill=MQWhite,
    inner sep=1.4pt
  },
  mqtt edge label/.style={
    font=\sffamily\scriptsize,
    text=MQTextMuted,
    fill=MQWhite,
    inner sep=1.4pt
  },
  mqtt topic label/.style={
    font=\ttfamily\scriptsize,
    text=MQText,
    fill=MQWhite,
    inner sep=1.4pt
  },
  mqtt route label/.style={
    font=\ttfamily\scriptsize,
    text=MQText,
    fill=MQWhite,
    inner sep=1.4pt
  },
  mqtt evidence runtime/.style={
    font=\sffamily\scriptsize\bfseries,
    text=MQSuccessLine,
    draw=MQEvidenceRuntimeLine,
    fill=MQEvidenceRuntimeFill,
    rounded corners=\MQRadiusXS,
    line width=.6pt,
    inner xsep=4pt,
    inner ysep=2pt
  },
  mqtt evidence source/.style={
    font=\sffamily\scriptsize\bfseries,
    text=MQControlLine,
    draw=MQEvidenceSourceLine,
    fill=MQEvidenceSourceFill,
    rounded corners=\MQRadiusXS,
    line width=.6pt,
    inner xsep=4pt,
    inner ysep=2pt
  }
}

% ---------------------------------------------------------------------------
% Canonical publication rules
% ---------------------------------------------------------------------------
%
% 1. Figure .tex + this style file are canonical technical sources.
% 2. Generated SVGs are never hand-edited.
% 3. SVG is the publication format for technical diagrams.
% 4. PNG is reserved for genuine runtime/UI evidence or raster-only input.
% 5. Desktop/mobile may share semantics but use independently art-directed geometry.
% 6. No arrow exists without real directional semantics.
% 7. Dashed arrows mean observation/secondary relation, not primary flow.
% 8. Containers denote real process/runtime/trust/configuration/lifecycle boundaries.
% 9. Color never carries the only semantic distinction.
% 10. Data plane, admin/control plane, trust boundaries and evidence classes use
%     the canonical semantic styles above.
% 11. Prefer explicit anchors/ports and orthogonal routing where clarity improves.
% 12. Preserve readable text size at actual GitHub desktop/mobile rendering widths.
% 13. Ordinary spacing follows the semantic spacing system above; avoid ad-hoc gaps.
% 14. Rounded corners follow the semantic radius scale above; avoid one-off radii.
% 15. Flat rendering is the default; elevation/shadows are explicit and semantically
%     justified.

% snodec-canonical-figure-system.tex
% MQTTSuite adapter for the authoritative SNode.C TikZ visual contract.
%
% Authority:
%   SNodeC/SNode.C-Book/assets/figures/src/snodec-figure-style.tex
%
% Load this AFTER mqttsystem-figure-system.tex. It preserves the MQTT-specific
% semantic vocabulary while overriding presentation constants and base styles
% with the canonical SNode.C palette, geometry, typography and connector grammar.

% Authoritative SNode.C palette.
\definecolor{snodecInk}{HTML}{17212B}
\definecolor{snodecMuted}{HTML}{4B5F68}
\definecolor{snodecRule}{HTML}{7A8B93}
\definecolor{snodecBlue}{HTML}{0B4F6C}
\definecolor{snodecBlueSoft}{HTML}{E3F2F8}
\definecolor{snodecGreen}{HTML}{1F7A68}
\definecolor{snodecGreenSoft}{HTML}{E5F4EF}
\definecolor{snodecOrange}{HTML}{C46A1B}
\definecolor{snodecOrangeSoft}{HTML}{FAEFE3}
\definecolor{snodecRed}{HTML}{B24A4A}
\definecolor{snodecRedSoft}{HTML}{F8E8E8}
\definecolor{snodecGraySoft}{HTML}{F1F3F2}
\colorlet{MQText}{snodecInk}
\colorlet{MQTextMuted}{snodecMuted}
\colorlet{MQLine}{snodecMuted!76}
\colorlet{MQBorder}{snodecRule!85}
\colorlet{MQSurface}{white}
\colorlet{MQSurfaceAlt}{snodecGraySoft}
\colorlet{MQDecisionFill}{snodecOrangeSoft}
\colorlet{MQDecisionLine}{snodecOrange!80}
\colorlet{MQSuccessFill}{snodecGreenSoft}
\colorlet{MQSuccessLine}{snodecGreen!78}
\colorlet{MQErrorFill}{snodecRedSoft}
\colorlet{MQErrorLine}{snodecRed!82}
\colorlet{MQControlFill}{snodecBlueSoft}
\colorlet{MQControlLine}{snodecBlue!78}
\colorlet{MQWarningFill}{snodecOrangeSoft}
\colorlet{MQWarningLine}{snodecOrange!80}
\colorlet{MQCrossLane}{snodecBlue!70}
\colorlet{MQDataPlaneFill}{snodecBlueSoft}
\colorlet{MQDataPlaneLine}{snodecBlue!78}
\colorlet{MQAdminPlaneFill}{snodecGraySoft}
\colorlet{MQAdminPlaneLine}{snodecRule!85}
\colorlet{MQTrustFill}{white}
\colorlet{MQTrustLine}{snodecRule!82}
\colorlet{MQEvidenceRuntimeFill}{snodecGreenSoft}
\colorlet{MQEvidenceRuntimeLine}{snodecGreen!78}
\colorlet{MQEvidenceSourceFill}{snodecBlueSoft}
\colorlet{MQEvidenceSourceLine}{snodecBlue!78}

% Authoritative geometry and spacing.
\renewcommand{\MQRadiusXS}{2.5pt}
\renewcommand{\MQRadiusS}{2.5pt}
\renewcommand{\MQRadiusM}{2.5pt}
\renewcommand{\MQRadiusL}{3pt}
\renewcommand{\MQRadiusXL}{3pt}
\renewcommand{\MQGapXS}{6mm}
\renewcommand{\MQGapS}{6mm}
\renewcommand{\MQGapM}{8mm}
\renewcommand{\MQGapL}{11mm}
\renewcommand{\MQGapXL}{14mm}
\renewcommand{\MQNodePadX}{3mm}
\renewcommand{\MQNodePadY}{2.5mm}
\renewcommand{\MQHorizontalGapTight}{6mm}
\renewcommand{\MQHorizontalGap}{8mm}
\renewcommand{\MQHorizontalGapLoose}{11mm}
\renewcommand{\MQHorizontalMilestoneGap}{14mm}
\renewcommand{\MQBranchGapTight}{6mm}
\renewcommand{\MQBranchGap}{8mm}
\renewcommand{\MQLaneGap}{14mm}
\renewcommand{\MQLanePadX}{4.5mm}
\renewcommand{\MQLanePadY}{4.5mm}
\renewcommand{\MQGroupPad}{4.5mm}
\renewcommand{\MQTrustPad}{4.5mm}
\renewcommand{\MQDesktopSafe}{8mm}
\renewcommand{\MQMobileSafe}{8mm}
\renewcommand{\MQMobileHorizontalGap}{8mm}
\renewcommand{\MQMobileLaneGap}{11mm}
\renewcommand{\MQDesktopMaxWidth}{160mm}
\renewcommand{\MQMobileMaxWidth}{100mm}

% Canonical desktop type scale. Mobile keeps the existing intentional uplift.
\renewcommand{\MQDesktopTypography}{%
  \renewcommand{\MQBodyFont}{\sffamily\small}%
  \renewcommand{\MQTitleFont}{\sffamily\bfseries\large}%
  \renewcommand{\MQSectionFont}{\sffamily\bfseries\small}%
  \renewcommand{\MQAnnotationFont}{\sffamily\footnotesize}%
  \renewcommand{\MQCodeFont}{\ttfamily\small}%
  \renewcommand{\MQCodeSmallFont}{\ttfamily\footnotesize}%
}

% Flat canonical node grammar. Keep this key list paragraph-free: pgfkeys does
% not accept \par tokens inside \tikzset.
\tikzset{
  mqtt elevated none/.style={},
  mqtt elevated subtle/.style={},
  mqtt elevated raised/.style={},
  mqtt process/.style={mqtt text,draw=snodecRule!85,fill=white,rounded corners=2.5pt,line width=.70pt,minimum height=10mm,minimum width=25mm,inner xsep=3mm,inner ysep=2.5mm},
  mqtt application/.style={mqtt text,draw=snodecGreen!78,fill=snodecGreenSoft,rounded corners=2.5pt,line width=.70pt,minimum height=10mm,minimum width=25mm,inner xsep=3mm,inner ysep=2.5mm},
  mqtt external/.style={mqtt process,fill=snodecGraySoft,draw=snodecRule!85},
  mqtt state/.style={mqtt process,fill=snodecBlueSoft,draw=snodecBlue!78},
  mqtt decision/.style={mqtt text,diamond,aspect=1.75,draw=snodecOrange!80,fill=snodecOrangeSoft,line width=.70pt,inner sep=1.2pt,minimum width=27mm,minimum height=14mm},
  mqtt success/.style={mqtt process,fill=snodecGreenSoft,draw=snodecGreen!78},
  mqtt error/.style={mqtt process,fill=snodecRedSoft,draw=snodecRed!82},
  mqtt control/.style={mqtt process,fill=snodecBlueSoft,draw=snodecBlue!78},
  mqtt warning/.style={mqtt process,fill=snodecOrangeSoft,draw=snodecOrange!80},
  mqtt note/.style={mqtt text,draw=snodecRule!85,fill=snodecGraySoft,rounded corners=2.5pt,line width=.58pt,font=\MQAnnotationFont,minimum height=8mm,inner xsep=3mm,inner ysep=2.5mm},
  mqtt config/.style={mqtt process,fill=snodecGraySoft,draw=snodecRule!85},
  mqtt data/.style={mqtt process,fill=white,draw=snodecRule!85},
  mqtt database/.style={mqtt text,cylinder,shape border rotate=90,aspect=.25,draw=snodecRule!85,fill=snodecGraySoft,line width=.70pt,minimum width=24mm,minimum height=13mm,inner xsep=3mm,inner ysep=2.5mm},
  mqtt figure frame/.style={draw=snodecRule!82,fill=white,rounded corners=3pt,line width=.58pt,inner xsep=4.5mm,inner ysep=4.5mm},
  mqtt lane/.style={draw=snodecRule!82,fill=none,dashed,rounded corners=3pt,line width=.58pt,inner xsep=4.5mm,inner ysep=4.5mm},
  mqtt group/.style={draw=snodecRule!82,fill=none,rounded corners=3pt,line width=.58pt,inner sep=4.5mm},
  mqtt data plane/.style={draw=snodecBlue!78,fill=snodecBlueSoft,rounded corners=3pt,line width=.58pt,inner sep=4.5mm},
  mqtt admin plane/.style={draw=snodecRule!85,fill=snodecGraySoft,rounded corners=3pt,line width=.58pt,inner sep=4.5mm},
  mqtt trust boundary/.style={draw=snodecRule!82,fill=none,dashed,rounded corners=3pt,line width=.58pt,inner sep=4.5mm}
}

% Authoritative connector grammar: round joins/caps and Latex tips.
\tikzset{
  mqtt connector/.style={line cap=round,line join=round},
  mqtt flow/.style={mqtt connector,-{Latex[length=2.2mm,width=1.6mm]},draw=snodecMuted!76,line width=.70pt},
  mqtt data flow/.style={mqtt connector,-{Latex[length=2.2mm,width=1.6mm]},draw=snodecBlue!70,line width=.70pt},
  mqtt control flow/.style={mqtt connector,-{Latex[length=2.2mm,width=1.6mm]},draw=snodecBlue!70,line width=.70pt},
  mqtt observation/.style={mqtt connector,-{Latex[length=2.0mm,width=1.4mm]},draw=snodecMuted!72,line width=.58pt,dashed},
  mqtt handoff/.style={mqtt connector,-{Latex[length=2.2mm,width=1.6mm]},draw=snodecBlue!70,line width=.70pt},
  mqtt association/.style={mqtt connector,draw=snodecMuted!68,line width=.58pt},
  mqtt contains/.style={mqtt connector,draw=snodecMuted!68,line width=.58pt},
  mqtt dependency/.style={mqtt connector,-{Latex[length=2.0mm,width=1.4mm]},draw=snodecMuted!68,line width=.58pt,densely dotted},
  mqtt section title/.style={mqtt no hyphenation,font=\MQSectionFont,text=snodecBlue,align=left,anchor=west},
  mqtt annotation/.style={mqtt no hyphenation,font=\MQAnnotationFont,text=snodecMuted,align=center},
  mqtt branch label/.style={mqtt no hyphenation,font=\MQAnnotationFont\bfseries,text=snodecInk,fill=white,inner sep=1pt},
  mqtt edge label/.style={mqtt no hyphenation,font=\MQAnnotationFont,text=snodecMuted,fill=white,inner sep=1pt},
  mqtt topic label/.style={mqtt no hyphenation,font=\MQCodeSmallFont,text=snodecInk,fill=white,inner sep=1pt}
}

\pgfdeclarelayer{framebackground}
\pgfsetlayers{framebackground,background,main}
\begin{document}
\begin{tikzpicture}[mqtt mobile]
\newcommand{\MQArrowRun}{8.3mm}
\tikzset{
  client/.style={mqtt external,minimum width=76mm,text width=68mm,minimum height=15mm},
  ingress/.style={mqtt process,minimum width=76mm,text width=68mm,minimum height=21mm},
  wssurface/.style={mqtt data,minimum width=35mm,text width=29mm,minimum height=27mm},
  adminsurface/.style={mqtt control,minimum width=35mm,text width=29mm,minimum height=27mm},
  riskcard/.style={mqtt warning,minimum width=76mm,text width=68mm,minimum height=30mm},
  deploymentcard/.style={mqtt note,minimum width=76mm,text width=68mm,minimum height=24mm}
}

% Equal sibling cards establish a common top edge. The router is placed exactly
% one canonical arrow run above that edge; midpoint routing makes both doglegs
% exactly symmetric without raw offsets or x/y shifts.
\node[wssurface] (ws)
  {\textbf{MQTT WebSocket}\\data route\\\texttt{/ · /ws · /mqtt}\\subprotocol \texttt{mqtt}};
\node[adminsurface,right=\MQHorizontalGapTight of ws] (admin)
  {\textbf{Operator routes}\\dashboard\\JSON \texttt{/api/mqtt/*}\\SSE};
\coordinate (surfacesTopMid) at ($(ws.north)!0.5!(admin.north)$);

\node[ingress,above=\MQArrowRun of surfacesTopMid] (router)
  {\textbf{shared Broker HTTP / HTTPS listener + Express router}\\browser, administration, SSE, and MQTT WebSocket share this ingress};
\coordinate (routerWsPort) at ($(router.south west)!.35!(router.south east)$);
\coordinate (routerAdminPort) at ($(router.south west)!.65!(router.south east)$);
\coordinate (wsMid) at ($(routerWsPort)!0.5!(ws.north)$);
\coordinate (adminMid) at ($(routerAdminPort)!0.5!(admin.north)$);
\draw[mqtt data flow]
  (routerWsPort) -- (routerWsPort |- wsMid) -| (ws.north);
\draw[mqtt control flow]
  (routerAdminPort) -- (routerAdminPort |- adminMid) -| (admin.north);

\node[client,above=\MQArrowRun of router] (client)
  {browser / HTTP client or MQTT-WebSocket client};
\draw[mqtt flow] (client.south) -- (router.north);

\begin{pgfonlayer}{background}
  \node[mqtt trust boundary,fit=(router)(ws)(admin)] (trust) {};
\end{pgfonlayer}
\node[mqtt annotation,text width=76mm,align=center,below=\MQLabelGap of trust] (trustLabel)
  {Dashed region: required external / trusted deployment boundary — not enforced by MQTTBroker};

\node[riskcard,below=\MQGapM of trustLabel] (risk)
  {\textbf{Current Broker exposure}\\No application authentication on operator routes. Wildcard CORS: \texttt{/api/mqtt/*} and \texttt{/api/mqtt/events}. Client-event JSON/logging can include the supplied MQTT password.};
\node[deploymentcard,below=\MQGapS of risk] (tls)
  {\textbf{Deployment control}\\HTTPS protects transport, not route authorization. Use trusted network policy or authenticated path-aware external access control.};

\node[inner sep=0pt,fit=(client)(trust)(trustLabel)(risk)(tls)] (contentBounds) {};
\coordinate (titleY) at ($(client.north)+(0,\MQTitleGap)$);
\node[mqtt section title,anchor=south west] (figureTitle)
  at (contentBounds.west |- titleY) {MQTTBroker HTTP exposure boundary};

\begin{pgfonlayer}{framebackground}
  \node[mqtt figure frame,fit=(figureTitle)(contentBounds)] {};
\end{pgfonlayer}
\end{tikzpicture}
\end{document}
